\subsection{Necesidad de seleccionar una biblioteca de compresión}

Una vez definido el simulador base, la siguiente decisión metodológica consistió en
seleccionar la biblioteca de compresión sin pérdida que permitiera materializar la
arquitectura propuesta sobre el vector de estado. En este trabajo, la compresión no se
plantea como una operación aislada de almacenamiento, sino como parte de un esquema
de gestión de memoria basado en particionamiento por bloques, descompresión selectiva
y reutilización temporal de datos. Por ello, la biblioteca elegida debía ser compatible con
un patrón de acceso repetido de lectura, modificación y escritura sobre fragmentos del
estado, y no limitarse a ofrecer una buena tasa de compresión en escenarios
secuenciales u orientados a archivo.

Bajo estas condiciones, la selección no debía responder únicamente al ratio de
compresión reportado por una biblioteca, sino a su adecuación como componente de una
ruta crítica de simulación numérica. En particular, interesaban bibliotecas utilizables
desde C/C++, con operación completamente sin pérdida, soporte eficiente para datos
residentes en memoria y comportamiento favorable dentro de una organización por
bloques independientes.
